\scanpage{22}
\begin{flushright}
\uline{tapis Quillen} (Notes du 6.9.68)
\end{flushright}

1.\quad Catégories formelles.

Catégorie formelle \inserted{$\mathcal C$} ``sur $X_0$'' i.e. dont \unsure{le 'objet|l'objet} des objets est le schéma $X_0$ (regardé comme \deleted{\unsure{ind-}} schéma formel i.e. \uline{ind-schéma} particulier).

Revient à la donnée d'un pro-Anneau strict \inserted{``admissible''} $A$ sur $X_0$, augmenté vers $\underline O_{X_0}${\reversemarginpar\margin{%
% Diagram in the left margin beside « augmenté vers ».
\diagram{\begin{tikzcd}[ampersand replacement=\&,row sep=large]
X_2 \arrow[d,shift right=4.4ex] \arrow[d,shift right=2.2ex] \arrow[d] \\
X_1 \arrow[u,shift right=2.3ex,bend right=9] \arrow[u,shift right=4.4ex,bend right=9] \arrow[d,shift right=2ex] \arrow[d] \\
X_0 \arrow[u,shift right=2.4ex,bend right=10]
\end{tikzcd}}
}}, l'idéal d'augmentation étant idéal de définition, avec deux structures de $\underline O_{X_0}$-modules (g. et dr.), et un hom. \inserted{de bi-$\underline O_{X_0}$-proanneaux} $A\xrightarrow{\Delta}A_{\mathrm{dr}}\otimes_{\underline O_{X_0}}{}_{\mathrm g}A\ \deleted{\unsure{\longrightarrow}}$, satisfaisant des axiomes évidents a) ``associativité'' de $\Delta$ b) unitarité g. et dr. vis à vis de l'augmentation. [Cela sera-t-il nécessairement \unsure{un|une} \deleted{catégorie} \uline{groupoïde} formel, i.e. y-aura-t-il une \uline{symétrie} $s$ de $A$, jouant vis-à-vis de $\Delta$ et de l'augmentation $\varepsilon$ le rôle \unsure{du|de} passage à l'inverse ?]

On pose $\Omega_{\mathcal C}={}$faisceau conormal de $X_0$ dans $X_1$ ; à priori c'est un pro-faisceau qu. ch. sur $X_0$, mais dans le cas où $A$ est $I$-\uline{adique} [$I$ l'idéal d'augmentation] c'est un module qu. ch. ordinaire.

On définit alors sur $\bigwedge^{*}\Omega$ une différentielle d'anneau gradué \inserted{$d:d_{\mathcal C}$} Sa valeur sur $\bigwedge^0\Omega=\underline O_{X_0}$ est
\[
df=\mathrm{pr}_2^{*}(f)-\mathrm{pr}_1^{\unsure{\alpha|*}}(f)\mod I^2
\]

% slanted
\margin{\uline{\unsure{n'interviendra|n'intervient|infinitésimale}} \uline{\unsure{pas|voir|vu} plus bas}}
\uline{Suite à développer} : $A=(A_\alpha)$, regardons $A_\alpha$ comme module pour la structure gauche, alors l'application $f\longmapsto\mathrm{pr}_2^{*}(f):\underline O_X\mathbin{-} A_\alpha$ est un op. diff. [d'ordre $n$ si $I_\alpha^{n+1}=0$], et par suite tte \deleted{\unsure{forme li}} section de $\underline{\mathrm{Hom}}_{\underline O_{\unsure{\alpha|X}}}(A_\alpha,\underline O_X)=\underline{\mathscr D}_{\alpha}^{\mathcal C}$ définit un op. différentiel dans $\underline O_X$. Alors $\mathscr D_{\unsure{0|\mathcal C}}^{\mathcal C}=\underline{\lim}_{\alpha}\mathscr D_{\alpha}^{\mathcal C}$ est un faisceau d'algèbres \inserted{\unsure{unitaires n.c.|unitaires ass.}} [grâce à l'application $\Delta$], qui s'envoie dans le faisceau des op différentiels [qui correspond

% archive sheet number: 21

